\clearpage
\begin{summarybox}
\textbf{\textcolor{CSummary}{一句话记住}}

\textbf{两个根号是两段路；把定点放在 $x$ 轴两侧，折线拉成直线时最短。}

\medskip
\textbf{\textcolor{CSummary}{它解决的题}}

两个根号相加，且同一个 $x$ 同时出现在两项中；分别求每项的最小值，往往不能同时取到。距离法可以一步找到整个和式的真正最小值。

\medskip
\textbf{\textcolor{CSummary}{一看就用的条件}}
\[
 f(x)=\sqrt{(x-m)^2+a^2}+\sqrt{(x-n)^2+b^2},
 \qquad a,b>0,\quad x\in\mathbb R.
\]

\medskip
\textbf{\textcolor{CSummary}{三步速答}}
\begin{description}[style=nextline,leftmargin=2.3em,labelsep=.5em,itemsep=.4em]
\item[\textbf{第 1 步：}] 先配方，认出 $m,n,a,b$。
\item[\textbf{第 2 步：}] 心里放三个点：$A(m,-a)$、$B(n,b)$、$P(x,0)$。
\item[\textbf{第 3 步：}] 直接写最短距离和交点位置：
\[
 \boxed{f_{\min}=\sqrt{(m-n)^2+(a+b)^2}},\qquad
 \boxed{x_0=\frac{bm+an}{a+b}}.
\]
\end{description}

完整值域为
\[
 \boxed{\left[\sqrt{(m-n)^2+(a+b)^2},+\infty\right)}.
\]

\medskip
\textbf{\textcolor{CSummary}{考试最容易丢分的三件事}}
\begin{itemize}[leftmargin=*,itemsep=.4em]
\item \textbf{不能分别取最小再相加}，除非两项刚好在同一个 $x$ 处最小。
\item \textbf{不能只写三点共线}；取等时 $P$ 必须落在线段 $AB$ 上。本题用上下两侧放点保证一定做到。
\item \textbf{定义域若受限，先检查 $x_0$}；落在允许范围外，就不能直接把完整实数范围的答案搬过来。
\end{itemize}

\medskip
\textbf{\textcolor{CSummary}{再多记一个漂亮的小结论}}

若 $m<n$，最短时
\[
 \boxed{(x_0-m):(n-x_0)=a:b}.
\]
也就是：\textbf{两个固定点离轴有多高，交点就按这个比例分配水平距离。}

\textbf{三种证明的选择：} 首选\textbf{画图拉直}（最直观）；想练代数就用\textbf{展开平方}（不需要求导）；遇到限制区间则借助\textbf{导数判断先减后增}（还能判断端点）。
\end{summarybox}
